%% P15 --- Jacobians, the chain rule and automatic differentiation
%%
%% Stub. The brief below is the contract for this program: what it must argue,
%% and what it must buy the reader. Writing the program means deleting the
%% \programstub{} block and replacing it with the program. If the brief turns
%% out to be wrong once the mathematics has been worked, change it and record
%% the contradiction in CLAUDE.md under *Resolved questions*.

\program{Jacobians, the chain rule and automatic differentiation}
\label{prog:P15}

\programstub{%
Argument: for vector-valued functions the chain rule multiplies Jacobians;
nobody forms a Jacobian; forward mode computes a Jacobian--vector product
and reverse mode a vector--Jacobian product, and which is cheaper depends
only on the shape of the problem. Payoff: what \code{loss.backward()}
actually computes, in full: why a scalar loss with many parameters makes
reverse mode cheaper by a factor of the parameter count; why reverse mode
must keep the activations, so memory scales with depth and batch; gradient
checkpointing as an explicit trade of recomputation for memory, with the
arithmetic done; why a numerical finite-difference gradient is a testing
tool and not an implementation; the three places autodiff silently gives you
something other than the derivative you meant (a non-differentiable point, a
detached tensor, an in-place write).

\medskip
\textbf{Planned length:} 65 frames.
}
